\documentclass[12pt]{article}
\usepackage[utf8]{inputenc}

\usepackage{lmodern}

\usepackage{enumitem}
\usepackage[margin=2cm]{geometry}

\usepackage{amsmath, amsfonts, amssymb}
\usepackage{graphicx}
%\usepackage{subfigure}
\usepackage{tikz}
\usepackage{pgfplots}
\usepackage{multicol}

\usepackage{comment}
\usepackage{url}
\usepackage{calc}
\usepackage{subcaption}
\usepackage[indent=0pt]{parskip}
\usepackage{animate}

\usepackage{array}
\usepackage{blkarray,booktabs, bigstrut}
\usepackage{bigints}

\pgfplotsset{compat=1.16}

% MATH commands
\newcommand{\ga}{\left\langle}
\newcommand{\da}{\right\rangle}
\newcommand{\oa}{\left\lbrace}
\newcommand{\fa}{\right\rbrace}
\newcommand{\oc}{\left[}
\newcommand{\fc}{\right]}
\newcommand{\op}{\left(}
\newcommand{\fp}{\right)}

\newcommand{\bi}{\mathbf{i}}
\newcommand{\bj}{\mathbf{j}}
\newcommand{\bk}{\mathbf{k}}
\newcommand{\bF}{\mathbf{F}}

\newcommand{\mR}{\mathbb{R}}

\newcommand{\ra}{\rightarrow}
\newcommand{\Ra}{\Rightarrow}

\newcommand{\sech}{\mathrm{sech}\,}
\newcommand{\csch}{\mathrm{csch}\,}
\newcommand{\curl}{\mathrm{curl}\,}
\newcommand{\dive}{\mathrm{div}\,}

\newcommand{\ve}{\varepsilon}
\newcommand{\spc}{\vspace*{0.5cm}}

\DeclareMathOperator{\Ran}{Ran}
\DeclareMathOperator{\Dom}{Dom}

\newcommand{\exo}[1]{\noindent\textcolor{red}{\fbox{\textbf{Problem {#1}}}\hrulefill}\\\\ }
\newcommand{\qu}[4]{\noindent\textcolor{#4}{\fbox{\textbf{Section {#1} | Problem {#2}}} \hrulefill{{\fbox{\textbf{{#3} Points}}}}\\}}

\newcommand{\semester}{Fall 2023}

\newcommand{\CVup}{%
\begin{tikzpicture}
\draw[black, <->, >=latex] (-0.33, 0.5) .. controls (-0.125, 0) and (0.125, 0) .. (0.33, 0.5);
\end{tikzpicture}}

\newcommand{\CVupInc}{%
\begin{tikzpicture}
\draw[black, ->, >=latex] (0,0) .. controls (0.2, 0) and (0.4, 0.2) .. (0.5, 0.5);
\end{tikzpicture}}

\newcommand{\CVupDec}{%
\begin{tikzpicture}[rotate=270]
\draw[black, ->, >=latex] (0,0) .. controls (0.2, 0) and (0.4, 0.2) .. (0.5, 0.5);
\end{tikzpicture}}

\newcommand{\CVdown}{%
\begin{tikzpicture}
\draw[black, <->, >=latex] (-0.33, -0.5) .. controls (-0.125, 0) and (0.125, 0) .. (0.33, -0.5);
\end{tikzpicture}}

\newcommand{\CVdownInc}{%
\begin{tikzpicture}
\draw[black, ->, >=latex] (-0.5, -0.5) .. controls (-0.5, -0.3) and (-0.5, -0.1) .. (0,0);
\end{tikzpicture}}

\newcommand{\CVdownDec}{%
\begin{tikzpicture}[rotate=-90]
\draw[black, ->, >=latex] (-0.5, -0.5) .. controls (-0.5, -0.3) and (-0.5, -0.1) .. (0,0);
\end{tikzpicture}}

\begin{document}
	\noindent \hrulefill \\
	MATH-244 \semester \hfill Practice Problems Solutions\\
	Section 15.7 \hfill Pierre-Olivier Paris{\'e} \\\vspace*{-1cm}
	
	\noindent\hrulefill
	
	\spc	

	\exo{6}
	\\
	The cylindrical coordinates are $(r, \theta , z)$ and the transformations to the cartesian coordinates $(x, y, z)$ are $x = r \cos \theta$, $y = r \sin \theta$, and $z = z$. So if $\theta = \pi/6$, then $x = r (\sqrt{2}/{2})$ and $y = r (1/2)$. So
		\begin{align*}
		2x/\sqrt{2} = r = 2y \quad \Rightarrow \quad x/\sqrt{2} = y .
		\end{align*}
	So the surface described by the equation $\theta = \pi/6$ is the plane $x/\sqrt{2} - y = 0$.
	
	\spc
	
	\exo{8}
	\\
	The transformation from cylindrical coordinates to the cartesian coordinates are $x = r \cos \theta$, $y = r \sin \theta$, and $z = z$. Since $\sin \theta = y/r$, the equation $r = 2 \sin \theta$ becomes
		\begin{align*}
		r = 2y/r \quad \Rightarrow \quad r^2 = 2y .
		\end{align*}
	But $r^2 = x^2 + y^2$, and so replacing in the last equation, we obtain
		\begin{align*}
		x^2 + y^2 = 2y \quad \iff \quad x^2 + (y - 1)^2 = 1 .
		\end{align*}
	This last equation represents a cylinder of radius 1 with center at $(0, 1)$.
	
	\spc
	
	\exo{12}
	\\
	The inequalities $0 \leq \theta \leq \pi/2$ means that we are in the first octant and the eigth octant. The inequalities $r \leq z \leq 2$ means that the value for $z$ is positive and it lies between the equations $z = r$ and $z = 2$. Since $r = \sqrt{x^2 + y^2}$, then $z$ lies above the cone $z = \sqrt{x^2 + y^2}$. Thus the solid is a quarter of a cone.
	
	\spc
	
	\exo{20}
	\\
	In cylindrical coordinate, the solid $E$ has the following description:
		\begin{align*}
		E = \{ (r, \theta , z) \, : \, 1 \leq r \leq 4 , \, 0 \leq \theta \leq 2\pi , \, 0 \leq z \leq y + 4 \} .
		\end{align*}
	So the triple integral can be rewritten as
		\begin{align*}
		\iiint_E (x - y) \, dV &= \int_0^{2\pi} \int_1^4 \int_0^{r \sin \theta +4} r \cos \theta - r \sin \theta \, dz r dr d \theta \\
		&= \int_0^{2\pi} \int_1^4 (\cos \theta - \sin \theta) (r \sin \theta + 4 ) r^2 \, dr d\theta \\
		&= \int_0^{2\pi} \int_1^4 r^3 \cos \theta \sin \theta + 4r^2 \cos \theta - r^3 \sin^2 \theta - 4 r^2 \sin \theta \, dr d\theta \\
		&= -255 \pi/4 .
		\end{align*}


\end{document}